\section*{Bab \ref{ch:b1:faults-folds-fabrics} --- Sesar, lipatan dan fabrik}

\begin{solution}{exo:b1:faults-folds-fabrics:1}
(a) \omterm{def:b1:faults-folds-fabrics:fault-types}{Sesar normal}. (b) \omterm{def:b1:faults-folds-fabrics:fault-types}{Sesar naik}, hampir murni \omterm{def:b1:geological-maps:attitude}{pergeseran-kemiringan}. (c) \omterm{def:b1:faults-folds-fabrics:fault-types}{Sesar geser-mendatar} \omterm{def:b1:faults-folds-fabrics:fault-types}{sinistral}. (d) \omterm{def:b1:faults-folds-fabrics:fault-types}{Sesar geser-mendatar} \omterm{def:b1:faults-folds-fabrics:fault-types}{dekstral}. (e) \omterm{def:b1:faults-folds-fabrics:fault-types}{Sesar geser-miring}: \omterm{def:b1:faults-folds-fabrics:fault-elements}{dinding gantung} telah bergerak turun dan sepanjang \omterm{def:b1:geological-maps:attitude}{jurus}, ke kiri; komponen geser-mendatarnya, $\cos\ang{40} = 0.77$ dari \omterm{def:b1:faults-folds-fabrics:fault-elements}{pergeseran bersih}, lebih besar daripada komponen \omterm{def:b1:geological-maps:attitude}{pergeseran-kemiringannya}, $0.64$, sehingga ia disebut sesar \omterm{def:b1:faults-folds-fabrics:fault-types}{sinistral} dengan komponen normal. (f) Dengan \omterm{def:b1:geological-maps:attitude}{kemiringan} \ang{25}, kurang dari \ang{45}, \omterm{def:b1:faults-folds-fabrics:fault-types}{sesar naik} itu adalah \omterm{def:b1:faults-folds-fabrics:fault-types}{sesar dorong}.
\end{solution}

\begin{solution}{exo:b1:faults-folds-fabrics:2}
Dengan $|\sin r| = 1$, lemparannya adalah $300\sin\ang{55} = \qty{245.7}{m}$ dan pergeseran horizontalnya $300\cos\ang{55} = \qty{172.1}{m}$.
\end{solution}

\begin{solution}{exo:b1:faults-folds-fabrics:3}
Untuk \omterm{def:b1:faults-folds-fabrics:fold-classes}{lipatan tegak} yang \omterm{def:b1:faults-folds-fabrics:fold-elements}{sayap-sayapnya} miring $\delta_1$ dan $\delta_2$ ke arah berlawanan, $\iota = \ang{180} - \delta_1 - \delta_2$. (a) \ang{110}: terbuka. (b) \ang{35}: ketat. (c) \ang{7}: isoklinal.
\end{solution}

\begin{solution}{exo:b1:faults-folds-fabrics:4}
\omterm{def:b1:faults-folds-fabrics:fault-types}{Sesar normal} miring $\ang{45} + \ang{35}/2 = \ang{62.5}$, \omterm{def:b1:faults-folds-fabrics:fault-types}{sesar dorong} $\ang{45} - \ang{17.5} = \ang{27.5}$. \omterm{def:b1:faults-folds-fabrics:fault-types}{Sesar geser-mendatar} konjugat masing-masing membentuk \ang{27.5} dengan $\sigma_1$, di kedua sisinya, sehingga keduanya bertemu pada \ang{55}.
\end{solution}

\begin{solution}{exo:b1:faults-folds-fabrics:5}
(a) Komponen geser-mendatar adalah $400\cos\ang{30} = \qty{346.4}{m}$ sepanjang \omterm{def:b1:geological-maps:attitude}{jurus}, komponen \omterm{def:b1:geological-maps:attitude}{pergeseran-kemiringan} $400\sin\ang{-30} = \qty{-200}{m}$, yaitu \qty{200}{m} ke bawah \omterm{def:b1:geological-maps:attitude}{kemiringan}; lemparannya $200\sin\ang{50} = \qty{153.2}{m}$ dan pergeseran horizontalnya $200\cos\ang{50} = \qty{128.6}{m}$. (b) \omterm{def:b1:geological-maps:attitude}{Jurusnya} $000$, sehingga $\mathbf s = (0, 1, 0)$;
$\mathbf u = (-\cos\ang{50}, 0, \sin\ang{50}) =
(-0.643, 0, 0.766)$. Maka
$\overrightarrow{AA'} = 400(0.866\,\mathbf s - 0.5\,\mathbf u) =
(128.6, 346.4, -153.2)$ dalam meter. Arahnya adalah $\arctan(128.6/346.4) =
\ang{20.4}$ dan \omterm{def:b1:field-lab-methods-1:trend-plunge}{penunjamannya} $\arcsin(153.2/400) = \ang{22.5}$, yang juga
$\arcsin(\sin\ang{30}\sin\ang{50})$. (c) \omterm{def:b1:faults-folds-fabrics:fault-types}{Sesar geser-miring}: \omterm{def:b1:faults-folds-fabrics:fault-elements}{rake} terletak
antara $0$ dan \ang{-90}, lebih dekat ke $0$: sesar \omterm{def:b1:faults-folds-fabrics:fault-types}{sinistral} dengan komponen normal.
\end{solution}

\begin{solution}{exo:b1:faults-folds-fabrics:6}
(a) Kedua sesar membentuk \ang{60} satu sama lain; $\sigma_1$ membagi dua sudut lancip, sehingga ia horizontal sepanjang $050$ (dan $230$); $\sigma_3$ horizontal sepanjang $140$;
$\sigma_2$ vertikal, sepanjang garis tempat kedua sesar vertikal bertemu. Setiap sesar membentuk \ang{30} dengan $\sigma_1$, sehingga $\ang{45} - \phi/2 =
\ang{30}$ dan $\phi = \ang{30}$. (b) Kedua sesar membagi tanah menjadi empat baji; baji-baji yang memuat $\sigma_1$, ke arah $050$ dan ke arah $230$, bergerak saling mendekat, dan yang memuat $\sigma_3$ bergerak saling menjauh, ke arah $140$ dan $320$. Ambil sesar yang berjurus $020$, di antara baji ke arah $050$ (di timurnya) dan baji ke arah $320$ (di baratnya): relatif terhadap baji barat, baji timur bergerak ke arah $230$ dan menjauh dari $320$, yang keduanya memiliki komponen ke arah $200$ sepanjang sesar. Seorang pengamat pada baji barat, menghadap timur, melihat baji timur bergerak ke arah $200$, ke kanan: sesar yang berjurus $020$ adalah \omterm{def:b1:faults-folds-fabrics:fault-types}{dekstral}. Argumen yang sama, atau simetri terhadap $\sigma_1$, membuat sesar yang berjurus $080$ \omterm{def:b1:faults-folds-fabrics:fault-types}{sinistral}.
\end{solution}

\begin{solution}{exo:b1:faults-folds-fabrics:7}
Normal ke atas adalah
\begin{align*}
  \mathbf n_1 &= (\sin\ang{40}\sin\ang{120}, \sin\ang{40}\cos\ang{120}, \cos\ang{40}) = (0.557, -0.321, 0.766),\\
  \mathbf n_2 &= (\sin\ang{30}\sin\ang{260}, \sin\ang{30}\cos\ang{260}, \cos\ang{30}) = (-0.492, -0.087, 0.866).
\end{align*}
Hasil kali silangnya adalah $(-0.211, -0.859, -0.207)$, yang menunjuk ke bawah: \omterm{def:b1:faults-folds-fabrics:fold-elements}{garis engsel} menunjam $\arcsin 0.207/0.909 = \ang{13.1}$ ke arah $\ang{180} + \arctan(0.211/0.859) =
\ang{193.8}$. $\mathbf n_1\cdot\mathbf n_2 = 0.417$, sudut antara normal-normalnya adalah \ang{65.3}, dan $\iota = \ang{114.7}$: \omterm{def:b1:faults-folds-fabrics:fold-classes}{lipatan terbuka}.
\end{solution}

\begin{solution}{exo:b1:faults-folds-fabrics:8}
$\mathbf n_1 = (0.321, 0.383, 0.866)$ untuk perlapisan dan
$\mathbf n_2 = (0.337, -0.925, 0.174)$ untuk kleavage. Hasil kali silang mereka, diambil ke bawah, memberi garis yang menunjam \ang{25.4} ke arah \ang{074.8}. Menurut
\cref{prop:b1:faults-folds-fabrics:intersection}, ini adalah arah \omterm{def:b1:faults-folds-fabrics:fold-elements}{garis engsel}: \omterm{def:b1:faults-folds-fabrics:fold-classes}{lipatan menunjam} sekitar \ang{25} ke arah timur--timur laut, yang seharusnya dikonfirmasi oleh \omterm{def:b1:geological-maps:geological-map}{singkapan} lain dari \omterm{def:b1:faults-folds-fabrics:fold-elements}{sayap} mana pun.
\end{solution}

\begin{solution}{exo:b1:faults-folds-fabrics:9}
(a) Pengamat di selatan sesar, menghadap utara, melihat bagian utara dike bergeser ke timur, ke kanan: sesarnya \omterm{def:b1:faults-folds-fabrics:fault-types}{dekstral}. (b) Dike adalah bidang vertikal yang tegak lurus sesar, sehingga pergeserannya sepanjang \omterm{def:b1:geological-maps:attitude}{jurus} adalah komponen geser-mendatar: \qty{350}{m}, yang merupakan seluruh pergeseran jika \omterm{def:b1:faults-folds-fabrics:fault-rocks}{gores garis} horizontal. (c) Pergeseran kemudian membentuk \ang{20} dengan \omterm{def:b1:geological-maps:attitude}{jurus} pada bidang sesar vertikal: komponen horizontalnya tetap \qty{350}{m}, sehingga \omterm{def:b1:faults-folds-fabrics:fault-elements}{pergeseran bersih} adalah $350/\cos\ang{20} = \qty{372.5}{m}$ dan komponen vertikalnya $350\tan\ang{20} = \qty{127.4}{m}$. Blok utara bergerak ke timur, dan \omterm{def:b1:faults-folds-fabrics:fault-rocks}{gores garis} turun ke timur, sehingga ia juga turun, sebesar \qty{127.4}{m}. Dike vertikal tidak dapat memperlihatkan gerakan vertikal itu, karena menggeser dike sepanjang dirinya sendiri tidak mengubah jejaknya.
\end{solution}

\begin{solution}{exo:b1:faults-folds-fabrics:10}
(a) Normal terhadap sesar membentuk $\ang{90} - \theta$ dengan $\sigma_1$, sehingga menurut \cref{prop:b1:stress-strain-rheology:mohr}, $\sigma_n' = \sigma_1'\sin^2\theta + \sigma_3'\cos^2\theta$ dan $\tau = (\sigma_1' - \sigma_3')\sin\theta\cos\theta$. Tergelincir dimulai ketika $\tau = \mu\sigma_n'$; dibagi dengan $\sigma_3'\cos^2\theta$ dan menuliskan $R = \sigma_1'/\sigma_3'$,
$(R - 1)\tan\theta = \mu(R\tan^2\theta + 1)$, sehingga $R\tan\theta(1 - \mu\tan\theta) = \tan\theta + \mu$, yaitu rumusnya.
(b) Dengan $\sigma_1$ vertikal, $\theta$ adalah \ang{90} dikurangi \omterm{def:b1:geological-maps:attitude}{kemiringan}: $\theta = \ang{30}$ memberi $R = 3.12$;
$\theta = \ang{45}$,
$R = 4.00$;
$\theta = \ang{55}$,
$R = 9.92$. (c) Penyebutnya lenyap ketika $\tan\theta = 1/\mu$,
$\theta = \ang{59.0}$; melampauinya, $R$ akan negatif: tidak ada rasio \omterm{def:b1:stress-strain-rheology:lithostatic}{tegangan efektif} positif yang dapat membuat sesar tergelincir, tidak peduli seberapa besar $\sigma_1'$ menjadi. \omterm{def:b1:faults-folds-fabrics:fault-types}{Sesar normal} yang miring kurang dari \ang{31.0} tidak dapat tergelincir di bawah $\sigma_1$ vertikal dengan $\mu = 0.6$. (d) Bidang yang paling baik orientasinya berada pada $\ang{45} - \phi/2$ dari $\sigma_1$. Untuk $\mu = 0.6$,
$\phi = \ang{31.0}$: \omterm{def:b1:faults-folds-fabrics:fault-types}{sesar normal} miring \ang{60.5}, \omterm{def:b1:faults-folds-fabrics:fault-types}{sesar dorong} \ang{29.5}. Untuk $\mu = 0.85$,
$\phi = \ang{40.4}$: \ang{65.2} dan \ang{24.8}.
\end{solution}

\begin{solution}{exo:b1:faults-folds-fabrics:11}
Melihat ke utara, timur ada di kanan, dan bidang yang miring ke barat naik ke arah kanan. (a) Di B perlapisan miring ke arah yang sama dengan kleavage, lebih curam: menurut \cref{prop:b1:faults-folds-fabrics:cleavage-bedding}(a), B berada pada \omterm{def:b1:faults-folds-fabrics:fold-elements}{sayap} terbalik; A dan C, tempat perlapisan kurang curam, tegak dengan benar. (b) \omterm{def:b1:faults-folds-fabrics:foliation}{Kleavage} naik pada \ang{60}; di A dan C perlapisan naik pada \ang{40}, dan rotasi lancip dari satu ke yang lain searah jarum jam, sebesar \ang{20}, menggerakkan bagian atas kleavage ke timur; di B perlapisan naik pada \ang{80} dan rotasinya berlawanan arah jarum jam, sebesar \ang{20}, menggerakkannya ke barat. (c) Arah berubah antara A dan B dan sekali lagi antara B dan C: ada \omterm{def:b1:faults-folds-fabrics:fold-elements}{engsel} di setiap interval. Di A, yang tegak dengan benar, bagian atas kleavage bergerak ke timur, menuju \omterm{def:b1:geological-maps:fold}{antiklin} di puncak \omterm{def:b1:faults-folds-fabrics:fold-elements}{sayap}: \omterm{def:b1:faults-folds-fabrics:fold-elements}{engsel} antara A dan B adalah \omterm{def:b1:geological-maps:fold}{antiklin}, dan yang antara B dan C adalah \omterm{def:b1:geological-maps:fold}{sinklin}. (d) Dalam notasi bukti, \omterm{def:b1:faults-folds-fabrics:fold-elements}{sayap-sayap} yang tegak dengan benar miring $\gamma - \iota/2 = \ang{40}$ dan \omterm{def:b1:faults-folds-fabrics:fold-elements}{sayap-sayap} terbalik $\gamma + \iota/2 =
\ang{80}$, sehingga $\gamma = \ang{60}$, seperti yang dikatakan kleavage, dan $\iota =
\ang{40}$: \omterm{def:b1:faults-folds-fabrics:fold-classes}{lipatan ketat}. \omterm{def:b1:faults-folds-fabrics:fold-elements}{Bidang aksial} berjurus utara dan miring \ang{60} ke barat, sikap $270/60$. \omterm{def:b1:faults-folds-fabrics:fold-elements}{Sayap-sayap} pendek telah berputar melewati vertikal dengan puncaknya ke timur: \omterm{def:b1:geological-maps:fold}{lipatan-lipatan} itu bervergensi ke timur, seperti dalam gambar bab ini.
\end{solution}

\begin{solution}{exo:b1:faults-folds-fabrics:12}
(a) Dalam \omterm{def:b1:geological-maps:fold}{lipatan} zigzag simetris setiap \omterm{def:b1:faults-folds-fabrics:fold-elements}{sayap} membentuk sudut $\iota/2$ dengan \omterm{def:b1:faults-folds-fabrics:fold-elements}{bidang aksial}, yang tegak lurus terhadap arah pemendekan; sebuah \omterm{def:b1:faults-folds-fabrics:fold-elements}{sayap} sepanjang $l$ karenanya membentang $l\sin(\iota/2)$ melintasi \omterm{def:b1:geological-maps:fold}{lipatan}, dan seluruh lapisan $L_0\sin(\iota/2)$. (b) $\sin\ang{30} = 0.5$: pemendekan sebesar \qty{50}{\percent}. (c) $\sin(\iota/2) = 0.7$ memberi $\iota = \ang{88.9}$.
(d) $\iota = \ang{156}$ dan $\sin\ang{78} = 0.978$: pemendekan sebesar \qty{2.2}{\percent}, yang juga $1 - \cos\ang{12}$.
\end{solution}

\begin{solution}{pb:b1:faults-folds-fabrics:1}
\textbf{1.} $\mathbf n = (\sin\delta\sin\alpha, \sin\delta\cos\alpha,
\cos\delta)$ memberi
\[
  \mathbf n_1 = (0.20855, 0.03415, 0.97742), \qquad
  \mathbf n_2 = (-0.20855, 0.03415, 0.97742).
\]

\textbf{2.} $\mathbf n_1\cdot\mathbf n_2 = -0.04349 + 0.00117 + 0.95535 =
0.91302$, dan sudut antara normal-normalnya adalah \ang{24.1}.

\textbf{3.} Menurut \cref{prop:b1:faults-folds-fabrics:hinge},
$\iota = \ang{180} - \ang{24.1} = \ang{155.9}$.

\textbf{4.} Antara \ang{170} dan \ang{90}: \omterm{def:b1:faults-folds-fabrics:fold-classes}{lipatan terbuka}.

\textbf{5.} \omterm{def:b1:faults-folds-fabrics:fold-elements}{Sayap-sayap} itu miring ke arah $080.7$ dan $279.3$, bukan tepat timur dan barat,
karena \omterm{def:b1:geological-maps:fold}{lipatannya} menunjam: garis-garis paling curamnya tidak berada dalam profil, yang hanya
menunjukkan \omterm{def:b1:cross-sections:apparent-dip}{kemiringan semu}. Dalam penampang timur--barat, menurut
\cref{thm:b1:cross-sections:apparent-dip}, $\tan\beta = \tan\ang{12.2}
\cos\ang{9.3} = 0.2134$, sehingga $\beta = \ang{12.0}$, dan $\ang{180} -
2\times\ang{12.0} = \ang{155.9}$ (profil itu miring \ang{2} dari penampang tersebut, yang tidak
mengubah apa pun pada ketelitian ini).

\textbf{6.} \omterm{def:b1:faults-folds-fabrics:fold-elements}{Sayapnya} sejajar dengan \omterm{def:b1:faults-folds-fabrics:fold-elements}{sayap} \omterm{def:b1:geological-maps:fold}{antiklin}, sehingga \omterm{def:b1:faults-folds-fabrics:fold-elements}{sudut antar-sayapnya} juga \ang{155.9}.

\textbf{7.} $\mathbf n_1\times\mathbf n_2 = (0.03415\times0.97742 -
0.97742\times0.03415,\; 0.97742\times(-0.20855) - 0.20855\times0.97742,\;
0.20855\times0.03415 + 0.03415\times0.20855) = (0, -0.40768, 0.01424)$.
Diambil ke bawah, $(0, 0.40768, -0.01424)$: \omterm{def:b1:faults-folds-fabrics:fold-elements}{garis engsel} berarah $000$ dan menunjam
$\arctan(0.01424/0.40768) = \ang{2.0}$ ke utara.

\textbf{8.} \omterm{def:b1:faults-folds-fabrics:fold-elements}{Bidang aksial} memuat \omterm{def:b1:faults-folds-fabrics:fold-elements}{garis engsel} dan membagi dua \omterm{def:b1:faults-folds-fabrics:fold-elements}{sayap}; normalnya terletak sepanjang
$\mathbf n_1 - \mathbf n_2 = (0.41710, 0, 0)$, horizontal dan timur--barat. Ia tegak dan berjurus
utara--selatan: $090/90$. \omterm{def:b1:geological-maps:fold}{Antiklin} itu tegak.

\textbf{9.} \omterm{def:b1:faults-folds-fabrics:fold-classes}{Lipatan terbuka} dan tegak, menunjam \ang{2.0} ke utara.

\textbf{10.} Menurut \cref{prop:b1:geological-maps:folds}, \omterm{def:b1:geological-maps:geological-map}{singkapan} \omterm{def:b1:geological-maps:fold}{antiklin} menutup ke arah
\omterm{def:b1:field-lab-methods-1:trend-plunge}{penunjaman}: ke utara. Intinya memuat batuan tertua: batupasir di dasar tutupan dan, di tempat sungai
paling dalam memotong, basemen sekis dan granit di bawahnya.

\textbf{11.} Sebagai \omterm{def:b1:faults-folds-fabrics:fold-elements}{bidang aksial}, kleavage tegak dan berjurus utara--selatan, $090/90$. Menurut
\cref{prop:b1:faults-folds-fabrics:intersection}, \omterm{def:b1:faults-folds-fabrics:lineation}{lineasi perpotongannya} dengan perlapisan sejajar
dengan \omterm{def:b1:faults-folds-fabrics:fold-elements}{garis engsel}: arah $000$, \omterm{def:b1:field-lab-methods-1:trend-plunge}{penunjaman} \ang{2.0}. Secara langsung: bidang $080.7/12.2$ bertemu
bidang vertikal yang berjurus utara sepanjang garis dengan \omterm{def:b1:field-lab-methods-1:trend-plunge}{penunjaman} $\arctan(\tan\ang{12.2}
\cos\ang{80.7}) = \ang{2.0}$ ke arah $000$.

\textbf{12.} Melihat ke utara, \omterm{def:b1:faults-folds-fabrics:fold-elements}{sayap} timur naik ke arah kiri pada \ang{12.0} dan kleavage vertikal:
rotasi lancip dari kleavage ke perlapisan, sebesar \ang{78.0}, berlawanan arah jarum jam dan
menggerakkan bagian atas kleavage ke barat. \omterm{def:b1:faults-folds-fabrics:fold-elements}{Sayap} itu tegak dengan benar, sehingga menurut
\cref{prop:b1:faults-folds-fabrics:cleavage-bedding}(b) \omterm{def:b1:faults-folds-fabrics:fold-elements}{engsel} \omterm{def:b1:geological-maps:fold}{antiklin} berada di barat, pada
$x = \qty{3.5}{km}$, seperti yang ditunjukkan peta.

\textbf{13.} \omterm{def:b1:geological-maps:attitude}{Jurusnya} $120 - 90 = 030$; bidang itu miring ke arah $120$, sehingga \omterm{def:b1:faults-folds-fabrics:fault-elements}{dinding gantung}
berada di sisi timur-tenggara dari jejak.

\textbf{14.} Sesar itu normal: $\sigma_1$ vertikal.

\textbf{15.} $\sigma_2$ horizontal sepanjang \omterm{def:b1:geological-maps:attitude}{jurus}, $030$--$210$;
$\sigma_3$ horizontal melintasinya,
$120$--$300$.

\textbf{16.} $\ang{45} + \phi/2 = \ang{60}$ memberi $\phi = \ang{30}$, dan
$\mu = \tan\ang{30} = 0.577$.

\textbf{17.} Himpunan konjugat juga berjurus $030$ dan miring \ang{60} ke arah sebaliknya:
$300/60$. Kedua himpunan bertemu pada $\ang{180} - 2\times\ang{60} = \ang{60}$, dengan $\sigma_1$
vertikal membagi dua sudut lancip.

\textbf{18.} Dengan $\mu = 0.85$,
$\phi = \ang{40.4}$ dan \omterm{def:b1:faults-folds-fabrics:fault-types}{sesar normal} yang paling baik orientasinya
miring \ang{65.2}; sesar Varden, pada \ang{60}, meleset \ang{5.2}. Menurut
\cref{exo:b1:faults-folds-fabrics:10}, ia tergelincir pada rasio
$\sigma_1'/\sigma_3' = (1 + 0.85\cot\ang{30})/(1 - 0.85\tan\ang{30}) = 4.85$,
dibandingkan dengan $4.68$ untuk bidang terbaik: ia hampir, tetapi tidak tepat, paling baik
orientasinya, dan akan tergelincir lagi pada rasio tegangan hanya \qty{4}{\percent} lebih tinggi.

\textbf{19.} Pergeserannya ke bawah \omterm{def:b1:geological-maps:attitude}{kemiringan}: $r = \ang{-90}$; \omterm{def:b1:faults-folds-fabrics:fault-elements}{vektor geser} berarah $120$ dan
menunjam \ang{60}.

\textbf{20.} $(\cos\ang{60}\sin\ang{120}, \cos\ang{60}\cos\ang{120},
-\sin\ang{60}) = (0.4330, -0.2500, -0.8660)$.

\textbf{21.} $e = 200/\tan\ang{60} = \qty{115.5}{m}$.

\textbf{22.} Timur: $115.47\sin\ang{120} = \qty{100.0}{m}$; utara:
$115.47\cos\ang{120} = \qty{-57.7}{m}$, yaitu \qty{57.7}{m} ke selatan.

\textbf{23.} $L = 200/\sin\ang{60} = \qty{230.9}{m}$, dan
$200^2 + 115.47^2 = 40\,000 + 13\,333 = 53\,333 = 230.94^2$. Lemparan dan pergeseran horizontal
adalah komponen vertikal dan horizontal pergeseran, yang terletak dalam bidang vertikal \omterm{def:b1:geological-maps:attitude}{kemiringan}:
keduanya adalah sisi segitiga siku-siku yang hipotenusanya adalah pergeseran.

\textbf{24.} Hingga meter terdekat, \omterm{def:b1:faults-folds-fabrics:fault-elements}{pergeseran bersih} pada sesar Varden adalah
\textbf{\qty[text-series-to-math=true]{231}{m}}.
\end{solution}
